\input{estilo}
% El contraste de tt-after: dos TT.now() que se solapan. No está en la
% pizarra; el texto lo dice en la misma frase que el caso que no se solapa.
\begin{document}
\begin{figura}\begin{tikzpicture}
  \newcommand\intervalo[6][0.5]{% [pos], y, x0, x1, rótulo, nombre
    \draw[intervalo, draw=segundo] (#3,#2) -- (#4,#2) node[pos=#1, above=3pt, etiq, text=segundo] (#6) {#5};
    \node[cod, text=segundo, below=4pt] at (#3,#2) {earliest};
    \node[cod, text=segundo, below=4pt] at (#4,#2) {latest};}
  % se solapan
  \intervalo{1.2}{0}{3.4}{\texttt{TT.now()} del evento 1}{f1}
  \intervalo[0.78]{0}{2.4}{5.8}{\texttt{TT.now()} del evento 2}{f2}
  \fill[acento, opacity=0.12] (2.4,1.05) rectangle (3.4,1.35);
  \fill[acento, opacity=0.12] (2.4,-0.15) rectangle (3.4,0.15);
  \draw[densely dotted, draw=acento, line width=0.5pt] (2.4,1.35) -- (2.4,0.15);
  % la guía del latest del evento 1 se corta donde está la palabra
  \node[cod, below=4pt] (l1) at (3.4,1.2) {\phantom{latest}};
  \draw[densely dotted, draw=acento, line width=0.5pt] (3.4,1.35) -- (l1.north) (l1.south) -- (3.4,0.15);
  \pic[scale=0.75] at (2.35,-1.15) {rechazo};
  \node[text=acento, anchor=west, align=left] at (2.6,-1.1)
    {se solapan: no hay seguridad de cuál fue primero\\[1pt]
     {\footnotesize \texttt{TT.after(latest\textsubscript{1})} es falso}};
  \tiempo{-0.6}{9.0}{-2.0}
\end{tikzpicture}\end{figura}
\end{document}
